\documentclass[10pt,a4paper]{amsart}
\usepackage{amsmath,amssymb,amsthm,graphicx}
% Retained preamble for the Stacks Project chapter snapshot in this directory.
% The chapter's own preamble.tex is a copy of the shared Stacks Project preamble; the block
% below is sliced MECHANICALLY, line for line, from it. Nothing was added, removed or
% redefined; the wrapper class, the xr-hyper/hyperref wiring and the externaldocument list
% are omitted because the assembler supplies the document class.
%
% lines taken from preamble.tex: 3, 5, 9, 11, 12, 156, 157, 158, ... (51 lines total)
%
% package coverage restored from the chapter's own preamble (2026-10-13): xy, xy-2cell

\usepackage[all]{xy}
\xyoption{2cell}
\UseAllTwocells

\newenvironment{reference}{\comment}{\endcomment}
\newenvironment{slogan}{\comment}{\endcomment}
\newenvironment{history}{\comment}{\endcomment}
\theoremstyle{plain}
\newtheorem{theorem}[subsection]{Theorem}
\newtheorem{proposition}[subsection]{Proposition}
\newtheorem{lemma}[subsection]{Lemma}
\theoremstyle{definition}
\newtheorem{definition}[subsection]{Definition}
\newtheorem{example}[subsection]{Example}
\newtheorem{exercise}[subsection]{Exercise}
\newtheorem{situation}[subsection]{Situation}
\theoremstyle{remark}
\newtheorem{remark}[subsection]{Remark}
\newtheorem{remarks}[subsection]{Remarks}
\numberwithin{equation}{subsection}
\def\lim{\mathop{\mathrm{lim}}\nolimits}
\def\colim{\mathop{\mathrm{colim}}\nolimits}
\def\Spec{\mathop{\mathrm{Spec}}}
\def\Hom{\mathop{\mathrm{Hom}}\nolimits}
\def\Ext{\mathop{\mathrm{Ext}}\nolimits}
\def\SheafHom{\mathop{\mathcal{H}\!\mathit{om}}\nolimits}
\def\SheafExt{\mathop{\mathcal{E}\!\mathit{xt}}\nolimits}
\def\Sch{\mathit{Sch}}
\def\Mor{\mathop{\mathrm{Mor}}\nolimits}
\def\Ob{\mathop{\mathrm{Ob}}\nolimits}
\def\Sh{\mathop{\mathit{Sh}}\nolimits}
\def\NL{\mathop{N\!L}\nolimits}
\def\CH{\mathop{\mathrm{CH}}\nolimits}
\def\proetale{{pro\text{-}\acute{e}tale}}
\def\etale{{\acute{e}tale}}
\def\QCoh{\mathit{QCoh}}
\def\Ker{\mathop{\mathrm{Ker}}}
\def\Im{\mathop{\mathrm{Im}}}
\def\Coker{\mathop{\mathrm{Coker}}}
\def\Coim{\mathop{\mathrm{Coim}}}
\DeclareMathSymbol{\boxtimes}{\mathbin}{AMSa}{"02}
\def\QCohstack{\mathcal{QC}\!\mathit{oh}}
\def\Cohstack{\mathcal{C}\!\mathit{oh}}
\def\Spacesstack{\mathcal{S}\!\mathit{paces}}
\def\Quotfunctor{\mathrm{Quot}}
\def\Hilbfunctor{\mathrm{Hilb}}
\def\Curvesstack{\mathcal{C}\!\mathit{urves}}
\def\Polarizedstack{\mathcal{P}\!\mathit{olarized}}
\def\Complexesstack{\mathcal{C}\!\mathit{omplexes}}
\def\Pic{\mathop{\mathrm{Pic}}\nolimits}
\def\Picardstack{\mathcal{P}\!\mathit{ic}}
\def\Picardfunctor{\mathrm{Pic}}
\def\Deformationcategory{\mathcal{D}\!\mathit{ef}}
\usepackage{cleveref}
\crefname{originallemma}{original lemma}{original lemmas}
\Crefname{originallemma}{Original lemma}{Original lemmas}
\makeatletter
\newlabel{curves-example-rational-tail}{{0E3H}{source}{}{sourcepointer-1}{}}
\newlabel{curves-example-rational-tail@cref}{{[originallemma][0][]0E3H}{[1][0][]source}{}{}{}}
\makeatother
\crefname{originallemma}{original lemma}{original lemmas}
\Crefname{originallemma}{Original lemma}{Original lemmas}
\makeatletter
\newlabel{curves-lemma-contracting-rational-tails}{{0E3L}{source}{}{sourcepointer-2}{}}
\newlabel{curves-lemma-contracting-rational-tails@cref}{{[originallemma][0][]0E3L}{[1][0][]source}{}{}{}}
\makeatother
\crefname{originallemma}{original lemma}{original lemmas}
\Crefname{originallemma}{Original lemma}{Original lemmas}
\makeatletter
\newlabel{curves-lemma-duality-dim-1-CM}{{0BS3}{source}{}{sourcepointer-3}{}}
\newlabel{curves-lemma-duality-dim-1-CM@cref}{{[originallemma][0][]0BS3}{[1][0][]source}{}{}{}}
\makeatother
\crefname{originallemma}{original lemma}{original lemmas}
\Crefname{originallemma}{Original lemma}{Original lemmas}
\makeatletter
\newlabel{curves-lemma-rational-tail-negative}{{0E63}{source}{}{sourcepointer-4}{}}
\newlabel{curves-lemma-rational-tail-negative@cref}{{[originallemma][0][]0E63}{[1][0][]source}{}{}{}}
\makeatother
\crefname{originallemma}{original lemma}{original lemmas}
\Crefname{originallemma}{Original lemma}{Original lemmas}
\makeatletter
\newlabel{curves-lemma-rr}{{0BS6}{source}{}{sourcepointer-5}{}}
\newlabel{curves-lemma-rr@cref}{{[originallemma][0][]0BS6}{[1][0][]source}{}{}{}}
\makeatother
\crefname{originallemma}{original lemma}{original lemmas}
\Crefname{originallemma}{Original lemma}{Original lemmas}
\makeatletter
\newlabel{lemma-genus}{{0E1J}{source}{}{sourcepointer-6}{}}
\newlabel{lemma-genus@cref}{{[originallemma][0][]0E1J}{[1][0][]source}{}{}{}}
\makeatother
\crefname{originallemma}{original lemma}{original lemmas}
\Crefname{originallemma}{Original lemma}{Original lemmas}
\makeatletter
\newlabel{section-relative-dualizing}{{0E6N}{source}{}{sourcepointer-7}{}}
\newlabel{section-relative-dualizing@cref}{{[originallemma][0][]0E6N}{[1][0][]source}{}{}{}}
\makeatother
\crefname{originallemma}{original lemma}{original lemmas}
\Crefname{originallemma}{Original lemma}{Original lemmas}
\makeatletter
\newlabel{spaces-duality-lemma-iso-on-RSheafHom}{{0E58}{source}{}{sourcepointer-8}{}}
\newlabel{spaces-duality-lemma-iso-on-RSheafHom@cref}{{[originallemma][0][]0E58}{[1][0][]source}{}{}{}}
\makeatother
\crefname{originallemma}{original lemma}{original lemmas}
\Crefname{originallemma}{Original lemma}{Original lemmas}
\makeatletter
\newlabel{spaces-duality-remark-iso-on-RSheafHom}{{0GG4}{source}{}{sourcepointer-9}{}}
\newlabel{spaces-duality-remark-iso-on-RSheafHom@cref}{{[originallemma][0][]0GG4}{[1][0][]source}{}{}{}}
\makeatother
\crefname{originallemma}{original lemma}{original lemmas}
\Crefname{originallemma}{Original lemma}{Original lemmas}
\makeatletter
\newlabel{spaces-perfect-lemma-flat-proper-perfect-direct-image-general}{{0CTM}{source}{}{sourcepointer-10}{}}
\newlabel{spaces-perfect-lemma-flat-proper-perfect-direct-image-general@cref}{{[originallemma][0][]0CTM}{[1][0][]source}{}{}{}}
\makeatother
\crefname{originallemma}{original lemma}{original lemmas}
\Crefname{originallemma}{Original lemma}{Original lemmas}
\makeatletter
\newlabel{spaces-perfect-lemma-open-where-cohomology-in-degree-i-rank-r-geometric}{{0E6A}{source}{}{sourcepointer-11}{}}
\newlabel{spaces-perfect-lemma-open-where-cohomology-in-degree-i-rank-r-geometric@cref}{{[originallemma][0][]0E6A}{[1][0][]source}{}{}{}}
\makeatother
\crefname{originallemma}{original lemma}{original lemmas}
\Crefname{originallemma}{Original lemma}{Original lemmas}
\makeatletter
\newlabel{varieties-lemma-check-invertible-sheaf-trivial}{{0B40}{source}{}{sourcepointer-12}{}}
\newlabel{varieties-lemma-check-invertible-sheaf-trivial@cref}{{[originallemma][0][]0B40}{[1][0][]source}{}{}{}}
\makeatother
\crefname{originallemma}{original lemma}{original lemmas}
\Crefname{originallemma}{Original lemma}{Original lemmas}
\makeatletter
\newlabel{varieties-lemma-no-sections-dual-nef}{{0E22}{source}{}{sourcepointer-13}{}}
\newlabel{varieties-lemma-no-sections-dual-nef@cref}{{[originallemma][0][]0E22}{[1][0][]source}{}{}{}}
\makeatother
\title{Stacks, Moduli of Curves: lemma-semistable}
\author{The Stacks Project Authors}
\date{stacks-project:moduli-curves.tex:local-snapshot-20261012; GNU FDL 1.2 or later}
\begin{document}
\maketitle
\noindent\textit{Editorial note.} Statement span is lines 2373-2384 of sources/stacks-moduli-curves/pinned-original-source.tex; proof span is lines 2386-2472, the complete author proof (85 body lines).\par\medskip
\noindent\textit{Editorial note.} Navigation stubs: 13 cited-but-not-included labels are supplied as non-load-bearing pointers with their pinned tags where the snapshot tag index has them.\par\medskip
\noindent\textit{Editorial note.} extra\_packages is empty on purpose: the retained preamble already loads xy with options and a second declaration clashes in the assembler wrapper.\par\medskip
\noindent\textit{Editorial note.} Text is retained verbatim, including the author's own typos and placeholders.\par\medskip
\noindent\textit{Editorial note.} Queue index 3 of queue\_moduli-curves.json (pb=4061, stmt 2373-2384, proof 2386-2472).\par\medskip
\noindent\textit{Allowed standard prerequisites (editorial).}\begin{itemize}
\item Explicit prerequisite boundary. The cited inputs are prerequisites; the argument in the retained span is not replaced by any cited result.
\end{itemize}\medskip
\ifdefined\hypertarget\hypertarget{sourcepointer-1}{}\fi
\noindent\textit{Source reference.} Original lemma 0E3H: sources/stacks-moduli-curves/pinned-original-source.tex, source lines 2472--2472. Cited by the retained statement or proof but not part of the collected spans; supplied only as a navigation pointer. The argument does not depend on its content.\par\medskip
\ifdefined\hypertarget\hypertarget{sourcepointer-2}{}\fi
\noindent\textit{Source reference.} Original lemma 0E3L: sources/stacks-moduli-curves/pinned-original-source.tex, source lines 2472--2472. Cited by the retained statement or proof but not part of the collected spans; supplied only as a navigation pointer. The argument does not depend on its content.\par\medskip
\ifdefined\hypertarget\hypertarget{sourcepointer-3}{}\fi
\noindent\textit{Source reference.} Original lemma 0BS3: sources/stacks-moduli-curves/pinned-original-source.tex, source lines 2472--2472. Cited by the retained statement or proof but not part of the collected spans; supplied only as a navigation pointer. The argument does not depend on its content.\par\medskip
\ifdefined\hypertarget\hypertarget{sourcepointer-4}{}\fi
\noindent\textit{Source reference.} Original lemma 0E63: sources/stacks-moduli-curves/pinned-original-source.tex, source lines 2472--2472. Cited by the retained statement or proof but not part of the collected spans; supplied only as a navigation pointer. The argument does not depend on its content.\par\medskip
\ifdefined\hypertarget\hypertarget{sourcepointer-5}{}\fi
\noindent\textit{Source reference.} Original lemma 0BS6: sources/stacks-moduli-curves/pinned-original-source.tex, source lines 2472--2472. Cited by the retained statement or proof but not part of the collected spans; supplied only as a navigation pointer. The argument does not depend on its content.\par\medskip
\ifdefined\hypertarget\hypertarget{sourcepointer-6}{}\fi
\noindent\textit{Source reference.} Original lemma 0E1J: sources/stacks-moduli-curves/pinned-original-source.tex, source lines 1128--1138. Cited by the retained statement or proof but not part of the collected spans; supplied only as a navigation pointer. The argument does not depend on its content.\par\medskip
\ifdefined\hypertarget\hypertarget{sourcepointer-7}{}\fi
\noindent\textit{Source reference.} Original lemma 0E6N: sources/stacks-moduli-curves/pinned-original-source.tex, source lines 2043--2126. Cited by the retained statement or proof but not part of the collected spans; supplied only as a navigation pointer. The argument does not depend on its content.\par\medskip
\ifdefined\hypertarget\hypertarget{sourcepointer-8}{}\fi
\noindent\textit{Source reference.} Original lemma 0E58: sources/stacks-moduli-curves/pinned-original-source.tex, source lines 2472--2472. Cited by the retained statement or proof but not part of the collected spans; supplied only as a navigation pointer. The argument does not depend on its content.\par\medskip
\ifdefined\hypertarget\hypertarget{sourcepointer-9}{}\fi
\noindent\textit{Source reference.} Original lemma 0GG4: sources/stacks-moduli-curves/pinned-original-source.tex, source lines 2472--2472. Cited by the retained statement or proof but not part of the collected spans; supplied only as a navigation pointer. The argument does not depend on its content.\par\medskip
\ifdefined\hypertarget\hypertarget{sourcepointer-10}{}\fi
\noindent\textit{Source reference.} Original lemma 0CTM: sources/stacks-moduli-curves/pinned-original-source.tex, source lines 2472--2472. Cited by the retained statement or proof but not part of the collected spans; supplied only as a navigation pointer. The argument does not depend on its content.\par\medskip
\ifdefined\hypertarget\hypertarget{sourcepointer-11}{}\fi
\noindent\textit{Source reference.} Original lemma 0E6A: sources/stacks-moduli-curves/pinned-original-source.tex, source lines 2472--2472. Cited by the retained statement or proof but not part of the collected spans; supplied only as a navigation pointer. The argument does not depend on its content.\par\medskip
\ifdefined\hypertarget\hypertarget{sourcepointer-12}{}\fi
\noindent\textit{Source reference.} Original lemma 0B40: sources/stacks-moduli-curves/pinned-original-source.tex, source lines 2472--2472. Cited by the retained statement or proof but not part of the collected spans; supplied only as a navigation pointer. The argument does not depend on its content.\par\medskip
\ifdefined\hypertarget\hypertarget{sourcepointer-13}{}\fi
\noindent\textit{Source reference.} Original lemma 0E22: sources/stacks-moduli-curves/pinned-original-source.tex, source lines 2472--2472. Cited by the retained statement or proof but not part of the collected spans; supplied only as a navigation pointer. The argument does not depend on its content.\par\medskip
\begin{lemma}
\label{lemma-semistable}
Let $f : X \to S$ be a prestable family of curves of genus $g \geq 1$.
Let $s \in S$ be a point of the base scheme. Let $m \geq 2$.
The following are equivalent
\begin{enumerate}
\item $X_s$ does not have a rational tail
(Algebraic Curves, Example \ref{curves-example-rational-tail}), and
\item $f^*f_*\omega_{X/S}^{\otimes m} \to \omega_{X/S}^{\otimes m}$,
is surjective over $f^{-1}(U)$ for some $s \in U \subset S$ open.
\end{enumerate}
\end{lemma}

\begin{proof}
Assume (2). Using the material in Section \ref{section-relative-dualizing}
we conclude that $\omega_{X_s}^{\otimes m}$ is
globally generated. However, if $C \subset X_s$
is a rational tail, then $\deg(\omega_{X_s}|_C) < 0$ by
Algebraic Curves, Lemma \ref{curves-lemma-rational-tail-negative}
hence $H^0(C, \omega_{X_s}|_C) = 0$ by
Varieties, Lemma \ref{varieties-lemma-check-invertible-sheaf-trivial}
which contradicts the fact that it is globally generated.
This proves (1).

\medskip\noindent
Assume (1). First assume that $g \geq 2$. Assumption (1) 
implies $\omega_{X_s}^{\otimes m}$ is globally generated,
see Algebraic Curves, Lemma \ref{curves-lemma-contracting-rational-tails}.
Moreover, we have
$$
\Hom_{\kappa(s)}(H^1(X_s, \omega_{X_s}^{\otimes m}), \kappa(s)) =
H^0(X_s, \omega_{X_s}^{\otimes 1 - m})
$$
by duality, see Algebraic Curves, Lemma \ref{curves-lemma-duality-dim-1-CM}.
Since $\omega_{X_s}^{\otimes m}$ is globally generated we find
that the restriction to each irreducible component has nonegative degree.
Hence the restriction of $\omega_{X_s}^{\otimes 1 - m}$ to each
irreducible component has nonpositive degree. Since
$\deg(\omega_{X_s}^{\otimes 1 - m}) = (1 - m)(2g - 2) < 0$ by Riemann-Roch
(Algebraic Curves, Lemma \ref{curves-lemma-rr}) we conclude that the $H^0$
is zero by Varieties, Lemma \ref{varieties-lemma-no-sections-dual-nef}.
By cohomology and base change we conclude that
$$
E = Rf_*\omega_{X/S}^{\otimes m}
$$
is a perfect complex whose formation commutes with arbitrary base change
(Derived Categories of Spaces, Lemma
\ref{spaces-perfect-lemma-flat-proper-perfect-direct-image-general}).
The vanishing proved above tells us that $E \otimes^\mathbf{L} \kappa(s)$
is equal to $H^0(X_s, \omega_{X_s}^{\otimes m})$ placed in degree $0$.
After shrinking $S$ we find $E = f_*\omega_{X/S}^{\otimes m}$
is a locally free $\mathcal{O}_S$-module placed in degree $0$
(and its formation commutes with arbitrary base change as
we've already said), see Derived Categories of Spaces, Lemma
\ref{spaces-perfect-lemma-open-where-cohomology-in-degree-i-rank-r-geometric}.
The map $f^*f_*\omega_{X/S}^{\otimes m} \to \omega_{X/S}^{\otimes m}$
is surjective after restricting to $X_s$. Thus it is surjective in
an open neighbourhood of $X_s$. Since $f$ is proper, this open
neighbourhood contains $f^{-1}(U)$ for some open neighbourhood
$U$ of $s$ in $S$.

\medskip\noindent
Assume (1) and $g = 1$. By
Algebraic Curves, Lemma \ref{curves-lemma-contracting-rational-tails}
the assumption (1) means that $\omega_{X_s}$ is isomorphic to
$\mathcal{O}_{X_s}$. If we can show that after shrinking $S$
the invertible sheaf $\omega_{X/S}$ because trivial, then
we are done. We may assume $S$ is affine. After shrinking $S$
further, we can write
$$
Rf_*\mathcal{O}_X = (\mathcal{O}_S \xrightarrow{0} \mathcal{O}_S)
$$
sitting in degrees $0$ and $1$
compatibly with further base change, see Lemma \ref{lemma-genus}.
By duality this means that
$$
Rf_*\omega_{X/S} = (\mathcal{O}_S \xrightarrow{0} \mathcal{O}_S)
$$
sitting in degrees $0$ and $1$\footnote{Use that
$Rf_*\omega_{X/S}^\bullet =
Rf_*R\SheafHom_{\mathcal{O}_X}(\mathcal{O}_X. \omega_{X/S}^\bullet) =
R\SheafHom_{\mathcal{O}_S}(Rf_*\mathcal{O}_X, \mathcal{O}_S)$
by Duality for Spaces, Lemma \ref{spaces-duality-lemma-iso-on-RSheafHom} and
Remark \ref{spaces-duality-remark-iso-on-RSheafHom}
and then that $\omega_{X/S}^\bullet = \omega_{X/S}[1]$ by
our definitions in Section \ref{section-relative-dualizing}.}.
In particular we obtain an isomorphism $\mathcal{O}_S \to f_*\omega_{X/S}$
which is compatible with base change since
formation of $Rf_*\omega_{X/S}$ is compatible with base change
(see reference given above).
By adjointness, we get a global section $\sigma \in \Gamma(X, \omega_{X/S})$.
The restriction of this section to the fibre $X_s$
is nonzero (a basis element in fact) and as
$\omega_{X_s}$ is trivial on the fibres,
this section is nonwhere zero on $X_s$.
Thus it nowhere zero in
an open neighbourhood of $X_s$. Since $f$ is proper, this open
neighbourhood contains $f^{-1}(U)$ for some open neighbourhood
$U$ of $s$ in $S$.
\end{proof}

\end{document}
